% Generated by manuscript/build_manuscript.py.
\section{Experiments}
\label{sec:experiments}


HoriScope is designed to replace separately optimized horizon-specific predictors with one prefix-consistent forecaster. The central empirical question is whether a single model can serve all requested horizons without sacrificing forecasting accuracy. We examine this question through benchmark comparisons and deployment costs, and then use controlled ablations and validation diagnostics to identify the contribution and behavior of the proposed components.

\subsection{Experimental setup}


\textbf{Datasets and metrics.} We evaluate HoriScope on seven public multivariate forecasting benchmarks: ETTm1, ETTm2, ETTh1, ETTh2, Weather, Electricity (ECL) and Solar. These benchmarks span electricity transformers, meteorology, electricity consumption and solar generation across different sampling frequencies and variable dimensions. Following the standard long-term forecasting protocol, we use prediction horizons $\mathcal H=\{96,192,336,720\}$ and report mean squared error (MSE) and mean absolute error (MAE), where lower values indicate better forecasts. Further metric definitions, dataset descriptions, statistics, data splits and preprocessing details are provided in \ref{app:experiment-details}.

\textbf{Model and baselines.} HoriScope is a decoder framework that can be coupled with different Encoders. We pair it with a lightweight patch-token MLP Encoder to isolate the contribution of the proposed decoder. The comparison covers four representative design families: (1) Transformer-based models, including SimpleTM, iTransformer and PatchTST \citep{chen2025simpletm,liu2024itransformer,nie2023patchtst}; (2) lightweight linear or attention-based forecasters and training frameworks, including TimeAlign, QDF and DLinear \citep{hu2026timealign,wang2026qdf,zeng2023dlinear}; (3) convolutional models, including TVNet and ModernTCN \citep{li2025tvnet,luo2024moderntcn}; and (4) multi-scale or decomposition-based models, including AMD and TimeMixer \citep{hu2025amd,wang2024timemixer}.

\textbf{Evaluation protocols.} The horizon-specific comparison tests whether one HoriScope model per dataset can retain forecasting accuracy against baselines optimized separately for each requested horizon. It therefore evaluates unified serving against the separate model specialization of the standard horizon-specific protocol. The complementary one-model-all-horizons comparison places every method under the same unified workflow. Each model is trained for the maximum horizon, and an $H$-step request is evaluated from its first $H$ outputs using the corresponding official fixed-$H$ test loader. This protocol evaluates how effectively each architecture serves all requested horizons through one model and one shared prediction trajectory.

\textbf{Implementation details.} Local experiments are implemented in Python 3.12.13 and PyTorch 2.9.0 with CUDA 12.8, and run on NVIDIA GeForce RTX 3090 GPUs. HoriScope is optimized under the BSCA objective using AdamW and a cosine learning-rate schedule \citep{loshchilov2019adamw,loshchilov2017sgdr}. All locally reproduced baselines are built from their official codebases, and any source-informed configuration adaptation is documented explicitly.

\subsection{Comparison with horizon-specific forecasters}
\label{sec:horizon-specific-comparison}


% Required packages: booktabs, graphicx, xcolor
\begin{table*}[t]
\centering
\scriptsize
\setlength{\tabcolsep}{2.0pt}
\resizebox{\textwidth}{!}{%
\begin{tabular}{l|cccccccccccccc}
\toprule
Dataset & HoriScope (Ours) & TimeAlign & QDF & AMD & SimpleTM & TVNet & iTransformer & TimeMixer & Leddam & ModernTCN & PatchTST & Crossformer & TimesNet & DLinear \\
\midrule
ETTm1 & \textcolor{red}{\textbf{0.329}}/\textcolor{red}{\textbf{0.363}} & \textcolor{blue}{\underline{0.341}}/\textcolor{blue}{\underline{0.368}} & 0.353/0.380 & 0.351/0.378 & 0.383/0.396 & 0.348/0.379 & 0.362/0.391 & 0.356/0.380 & 0.360/0.390 & 0.351/0.381 & 0.353/0.376 & 0.416/0.435 & 0.399/0.406 & 0.356/0.378 \\
ETTm2 & \textcolor{blue}{\underline{0.248}}/\textcolor{red}{\textbf{0.304}} & \textcolor{red}{\textbf{0.245}}/\textcolor{red}{\textbf{0.304}} & 0.257/0.315 & 0.254/0.315 & 0.280/0.325 & 0.251/\textcolor{blue}{\underline{0.311}} & 0.269/0.329 & 0.257/0.318 & 0.265/0.321 & 0.253/0.314 & 0.256/0.317 & 0.518/0.501 & 0.292/0.331 & 0.259/0.324 \\
ETTh1 & \textcolor{red}{\textbf{0.390}}/\textcolor{red}{\textbf{0.415}} & 0.418/0.431 & 0.439/0.442 & 0.412/0.428 & 0.428/0.433 & 0.407/0.421 & 0.439/0.448 & 0.427/0.441 & 0.415/0.431 & \textcolor{blue}{\underline{0.404}}/\textcolor{blue}{\underline{0.420}} & 0.419/0.437 & 0.440/0.463 & 0.488/0.452 & 0.425/0.439 \\
ETTh2 & \textcolor{red}{\textbf{0.306}}/\textcolor{red}{\textbf{0.364}} & 0.347/0.387 & 0.353/0.394 & 0.366/0.406 & 0.356/0.392 & 0.324/\textcolor{blue}{\underline{0.377}} & 0.374/0.406 & 0.349/0.397 & 0.346/0.392 & \textcolor{blue}{\underline{0.322}}/0.379 & 0.351/0.395 & 0.809/0.658 & 0.390/0.417 & 0.431/0.447 \\
Weather & \textcolor{red}{\textbf{0.214}}/\textcolor{red}{\textbf{0.245}} & \textcolor{blue}{\underline{0.218}}/\textcolor{blue}{\underline{0.247}} & 0.237/0.271 & 0.225/0.265 & 0.243/0.271 & 0.221/0.261 & 0.233/0.271 & 0.226/0.264 & 0.227/0.264 & 0.224/0.264 & 0.226/0.264 & 0.228/0.287 & 0.259/0.287 & 0.242/0.293 \\
ECL & \textcolor{red}{\textbf{0.151}}/\textcolor{red}{\textbf{0.244}} & \textcolor{blue}{\underline{0.156}}/\textcolor{blue}{\underline{0.246}} & 0.167/0.260 & 0.162/0.257 & 0.168/0.262 & 0.165/0.254 & 0.164/0.261 & 0.173/0.272 & 0.162/0.257 & \textcolor{blue}{\underline{0.156}}/0.253 & 0.159/0.253 & 0.181/0.279 & 0.193/0.295 & 0.166/0.264 \\
Solar & \textcolor{red}{\textbf{0.188}}/\textcolor{red}{\textbf{0.208}} & 0.196/\textcolor{blue}{\underline{0.217}} & 0.208/0.258 & 0.205/0.248 & \textcolor{blue}{\underline{0.191}}/0.251 & 0.228/0.277 & 0.202/0.248 & 0.193/0.252 & 0.224/0.265 & 0.228/0.282 & 0.194/0.245 & 0.205/0.232 & 0.244/0.334 & 0.223/0.226 \\
\midrule
Average & \textcolor{red}{\textbf{0.261}}/\textcolor{red}{\textbf{0.306}} & \textcolor{blue}{\underline{0.274}}/\textcolor{blue}{\underline{0.314}} & 0.288/0.331 & 0.282/0.328 & 0.293/0.333 & 0.278/0.326 & 0.292/0.336 & 0.283/0.332 & 0.286/0.331 & 0.277/0.327 & 0.280/0.326 & 0.399/0.408 & 0.324/0.360 & 0.300/0.339 \\
\bottomrule
\end{tabular}%
}
\caption{Comparison with horizon-specific forecasters. Each dataset entry reports MSE/MAE averaged over $H\in\{96,192,336,720\}$, and the final row is the unweighted mean over seven datasets. HoriScope uses one unified model per dataset, whereas each baseline follows its horizon-specific protocol. Full horizon-wise results, source roles and protocol details are provided in ~\ref{app:full-results}. Best and second-best displayed values are marked in red bold and blue underline, respectively.}
\label{tab:main_horiscope}
\end{table*}



Table~\ref{tab:main_horiscope} compares one HoriScope model per dataset with baselines optimized separately for each requested horizon and reports the four-horizon mean for each dataset. Complete horizon-wise results are provided in \ref{app:full-results}. HoriScope obtains the best result in 13 of the 14 dataset--metric comparisons and the second-best result in the remaining comparison. Compared with TimeAlign, the strongest aggregate horizon-specific baseline, HoriScope reduces the seven-dataset average MSE and MAE by 4.94\% and 2.54\%, respectively. The advantage persists across datasets with markedly different dimensionalities. Over the four low-dimensional ETT benchmarks, HoriScope improves the average MSE and MAE over TimeAlign by 5.69\% and 2.89\%; on the high-dimensional ECL and Solar datasets, the corresponding improvements are 3.90\% and 2.27\%. These results use a lightweight patch-token MLP Encoder, indicating that HoriScope can construct accurate future trajectories from comparatively shallow history representations. Overall, HoriScope improves forecasting accuracy over the evaluated horizon-specific baselines while consolidating the four requested horizons into one model.

\subsection{One-model-all-horizons evaluation}
\label{sec:one-model-comparison}


Section 5.2 shows that one unified HoriScope model can outperform baselines trained separately for each of the four horizons. A stricter comparison places every method under the same one-model-all-horizons workflow. For each dataset, a single model produces the maximum-length forecast, and requests for $H\in\{96,192,336\}$ are evaluated from the corresponding output prefixes without retraining or reselection. This setting maintains consistency across overlapping predictions within each method and provides a more deployment-relevant comparison without horizon-specific model specialization.

% Required packages: booktabs, graphicx, xcolor
\begin{table*}[t]
\centering
\scriptsize
\setlength{\tabcolsep}{2.0pt}
\resizebox{\textwidth}{!}{%
\begin{tabular}{l|cccccccc}
\toprule
Dataset & ISCF-BSCA (Ours) & TimeAlign & QDF & AMD & SimpleTM & iTransformer & PatchTST & DLinear \\
\midrule
ETTm1 & \textcolor{red}{\textbf{0.329}}/\textcolor{red}{\textbf{0.363}} & \textcolor{blue}{\underline{0.345}}/\textcolor{blue}{\underline{0.371}} & 0.352/0.384 & 0.354/0.380 & 0.384/0.399 & 0.409/0.415 & 0.354/0.384 & 0.359/0.381 \\
ETTm2 & \textcolor{red}{\textbf{0.248}}/\textcolor{red}{\textbf{0.304}} & \textcolor{blue}{\underline{0.250}}/\textcolor{blue}{\underline{0.308}} & 0.257/0.316 & 0.255/0.316 & 0.285/0.330 & 0.294/0.339 & 0.258/0.316 & 0.292/0.359 \\
ETTh1 & \textcolor{red}{\textbf{0.390}}/\textcolor{red}{\textbf{0.415}} & 0.420/0.433 & 0.452/0.453 & \textcolor{blue}{\underline{0.416}}/\textcolor{blue}{\underline{0.430}} & 0.434/0.436 & 0.468/0.458 & 0.429/0.440 & 0.474/0.480 \\
ETTh2 & \textcolor{red}{\textbf{0.306}}/\textcolor{red}{\textbf{0.364}} & 0.368/0.399 & \textcolor{blue}{\underline{0.343}}/\textcolor{blue}{\underline{0.392}} & 0.370/0.411 & 0.356/0.393 & 0.385/0.408 & 0.352/0.395 & 0.415/0.439 \\
Weather & \textcolor{red}{\textbf{0.214}}/\textcolor{red}{\textbf{0.245}} & \textcolor{blue}{\underline{0.220}}/\textcolor{blue}{\underline{0.250}} & 0.260/0.295 & 0.226/0.266 & 0.254/0.280 & 0.262/0.284 & 0.231/0.268 & 0.247/0.301 \\
ECL & \textcolor{red}{\textbf{0.151}}/\textcolor{red}{\textbf{0.244}} & \textcolor{blue}{\underline{0.157}}/\textcolor{blue}{\underline{0.246}} & 0.183/0.281 & 0.163/0.257 & 0.169/0.265 & 0.177/0.269 & 0.161/0.255 & 0.167/0.265 \\
Solar & \textcolor{red}{\textbf{0.188}}/\textcolor{red}{\textbf{0.208}} & 0.193/\textcolor{blue}{\underline{0.219}} & 0.207/0.268 & 0.204/0.250 & \textcolor{blue}{\underline{0.190}}/0.247 & 0.229/0.261 & 0.194/0.247 & 0.253/0.313 \\
\midrule
Average & \textcolor{red}{\textbf{0.261}}/\textcolor{red}{\textbf{0.306}} & \textcolor{blue}{\underline{0.279}}/\textcolor{blue}{\underline{0.318}} & 0.293/0.341 & 0.284/0.330 & 0.296/0.336 & 0.318/0.348 & 0.283/0.329 & 0.315/0.363 \\
\bottomrule
\end{tabular}%
}
\caption{One-model-all-horizons comparison. Each method uses one unified model per dataset, and each entry reports MSE/MAE averaged over $H\in\{96,192,336,720\}$. The final row is the unweighted mean over seven datasets. Full horizon-wise results and protocol details are provided in ~\ref{app:full-results}. Best and second-best displayed values are marked in red bold and blue underline, respectively.}
\label{tab:main_iscf_bsca_one_model}
\end{table*}



Table~\ref{tab:main_horiscope_one_model} summarizes the four-horizon mean for each dataset, with complete horizon-wise results provided in \ref{app:full-results}. Under the shared one-model constraint, HoriScope ranks first in all 14 dataset--metric comparisons. Averaged over all seven datasets and four horizons, it reduces MSE and MAE by 6.45\% and 3.72\% relative to TimeAlign, the next strongest complete baseline in Table~\ref{tab:main_horiscope_one_model}. These gains exceed the corresponding 4.94\% and 2.54\% improvements in the horizon-specific comparison by 1.51 and 1.18 percentage points, respectively, showing that the advantage of HoriScope becomes more pronounced when every method follows the same UVHF workflow.

These results establish HoriScope as an effective forecaster for UVHF. Its scope-indexed forecast field constructs one prefix-consistent trajectory, while Target-Adaptive Scope Allocation integrates predictive information across different sharing granularities for each future step. The consistently strong results across the four evaluated horizons support this output-side design for forecasting requests ranging from short to long endpoints.

\subsection{Accuracy and system cost}
\label{sec:system-cost}


Forecast accuracy is only one consideration in practical forecasting services. Runtime memory and model storage also determine deployment and management costs, particularly in resource-constrained settings. We therefore measure the peak allocated GPU memory of the complete inference service and the total checkpoint storage required to serve all four horizons. Fig.~\ref{fig:accuracy-system-cost} compares forecasting accuracy with these two costs.



\begin{figure*}[t]
\centering
\includegraphics[width=0.90\textwidth]{figures/figure_05_accuracy_system_cost.pdf}
\caption{Accuracy and system cost for four-horizon forecasting. Macro MSE is plotted against total checkpoint storage, and bubble area represents peak allocated inference memory. The storage axis is logarithmic; lower-left positions indicate a more favorable trade-off.}
\label{fig:accuracy-system-cost}
\end{figure*}


As shown in Fig.~\ref{fig:accuracy-system-cost}, HoriScope attains the lowest macro MSE among the eight displayed methods, at 0.261, while requiring 17.677 MiB of checkpoint storage and 38.817 MiB of peak allocated inference memory. Compared with TimeAlign, AMD, iTransformer, PatchTST and TimeMixer, the unified model reduces checkpoint storage by 30.12--92.01\% and peak memory by 18.01--83.69\%. Relative to TimeAlign in particular, HoriScope lowers MSE by 4.94\%, checkpoint storage by 81.48\% and peak memory by 64.42\%.

\subsection{Component and training-objective ablations}
\label{sec:ablation}


To determine whether the individual HoriScope components contribute to forecasting performance, we conduct strictly controlled ablations in which each variant is trained end to end under the same protocol. \textbf{w/o BSCA} retains the Uniform-Prefix Forecasting Loss but removes the Scope-Wise Forecasting Loss and Allocation-Balance Regularizer. \textbf{w/o Target-Adaptive Allocation} replaces learned Scope Probabilities with equal non-adaptive fusion. \textbf{Shared Scope Projection} replaces the scope-specific history projections with one shared projection, whereas \textbf{Fixed Scope ($s=144$)} retains only the preregistered middle scope.

\input{tables/table_3.tex}


As reported in Table~\ref{tab:core-ablation}, Full HoriScope achieves a five-dataset macro MSE/MAE of 0.305/0.344 and ranks first in all 12 dataset and average metric columns. Removing BSCA causes the largest aggregate decline, while equal fusion, Shared Scope Projection and Fixed Scope also increase both error metrics. Full HoriScope improves MSE and MAE over every control on all five datasets.

The controlled comparisons support the utility of BSCA, Target-Adaptive Scope Allocation, scope-specific projections and multi-scope generation. Their joint design yields the strongest aggregate forecasting performance.

\subsection{Scope diversity and allocation behavior}
\label{sec:scope-behavior}


The ablation results support the contribution of the individual HoriScope components to forecasting accuracy. We further examine whether the complete model exhibits the intended internal behavior. The analysis asks whether jointly trained scope lines retain distinct forecast signals and whether the Target-Adaptive Allocation Path produces diverse preferences across future regions. Region-specific forecast generation requires both properties: scope lines must avoid trajectory collapse, and allocation must vary across the future domain.

We concatenate the region-wise predictions from each scope line into its complete Scope-conditioned Forecast. Direct overlay obscures the differences among the five trajectories, so Fig.~\ref{fig:scope-allocation-behavior}b visualizes their deviations from the fused trajectory, $\mathcal F_\tau^{(s)}-\widehat y_\tau$. The profiles exhibit distinct temporal structures, with a mean pairwise absolute disagreement of 0.0812 on the normalized scale. The scope-indexed forecast field therefore retains distinct candidate signals for region-wise integration and avoids collapse to repeated trajectories.

Fig.~\ref{fig:scope-allocation-behavior}a shows how the Target-Adaptive Allocation Path combines these heterogeneous scope signals. The fused maximum-horizon forecast is coloured by the scope receiving the highest mean Scope Probability within each future region. All five scopes attain the highest weight in at least one of the eight regions, indicating that allocation avoids collapse to a fixed regional preference. Fig.~\ref{fig:scope-allocation-behavior}c displays the complete per-step Scope Probabilities, whose relative ordering varies across future steps and regions. The farthest region assigns its highest mean weight to the largest sharing scope, $s=720$, consistent with the role of broader information sharing in distal forecasting. The three panels illustrate how distinct scope signals and region-dependent soft allocation support region-specific forecast generation. These signals are fused into one shared trajectory, while prefix-based inference preserves CHPC with a maximum absolute CHPD of 0 across all 20 dataset--horizon checks.



\begin{figure*}[t]
\centering
\includegraphics[width=1.00\textwidth]{figures/figure_06_scope_allocation_behavior.pdf}
\caption{Scope diversity and sample-specific allocation behavior in HoriScope. \textbf{a}, The fused forecast; line colour and the lower strip identify the highest-weight scope after averaging Scope Probabilities within each future region. \textbf{b}, Deviations of the five Scope-conditioned Forecasts from the fused forecast on a shared scale. \textbf{c}, Per-step Scope Probabilities; white lines mark future-region boundaries.}
\label{fig:scope-allocation-behavior}
\end{figure*}


\subsection{Generalization studies}
\label{sec:generalization}


HoriScope is designed as an Encoder-agnostic plug-in decoder that operates on Encoder-produced history representations. We examine this versatility using DLinear and PatchTST as representatives of lightweight linear and Transformer-based Encoders, respectively \citep{zeng2023dlinear,nie2023patchtst}. For each backbone, we replace the Original Decoder with HoriScope and train the resulting model end to end using BSCA.



\begin{figure*}[t]
\centering
\includegraphics[width=0.90\textwidth]{figures/figure_07_decoder_transfer.pdf}
\caption{Generalization across two forecasting backbones. Four-horizon mean MSE is reported for \textbf{a}, DLinear and \textbf{b}, PatchTST with either the Original Decoder or HoriScope.}
\label{fig:decoder-transfer}
\end{figure*}


As shown in Fig.~\ref{fig:decoder-transfer}, replacing the Original Decoder with HoriScope improves forecasting performance on both example backbones. For DLinear, the replacement lowers MSE on all three datasets and reduces the macro average from 0.303 to 0.286, corresponding to an improvement of 5.61\%. PatchTST shows smaller but consistent gains, with macro MSE decreasing from 0.282 to 0.276, corresponding to an improvement of 2.13\%. Improvements with both a lightweight linear backbone and a Transformer-based backbone support the intended role of HoriScope as a general plug-in decoder for UVHF.
